\midsectiontitle{Atributos do Universo}

Os atributos definem as capacidades essenciais de um personagem, moldando sua força física, agilidade, resistência, percepção, intelecto e conexão com o divino.  
Cada atributo influencia diretamente testes, combates e interações narrativas.

\noindent

% =========================
% MIG
% =========================

\originlink{att:might}
\noindent
\textbf{\baskerfamily\large MIG — Might (Poder)}\\
A força bruta do corpo e da vontade física. MIG governa o dano em combate corpo a corpo, a capacidade de empunhar armas pesadas e realizar feitos de pura potência, como arrombar portões, erguer estruturas ou subjugar inimigos pela força.


% =========================
% GRA
% =========================
\originlink{att:grace}
\noindent
\textbf{\baskerfamily\large GRA — Grace (Graça)}\\
Expressa velocidade, precisão e harmonia de movimento. GRA afeta esquivas, acrobacias, ataques rápidos e qualquer ação que exija controle corporal refinado. Personagens com alta Graça parecem mover-se como se guiados por mãos divinas.

% =========================
% FOR
% =========================
\originlink{att:fortitude}
\noindent
\textbf{\baskerfamily\large FOR — Fortitude (Fortitude)}\\
A resistência do corpo e do espírito. FOR define pontos de vida, resistência a doenças, venenos e efeitos prolongados. Um personagem com alta Fortitude sobrevive onde outros tombam.

% =========================
% WIS
% =========================
\originlink{att:wisdom}
\noindent
\textbf{\baskerfamily\large WIS — Wisdom (Sabedoria)}\\
Representa conhecimento, discernimento e energia arcana. WIS controla o aprendizado, o uso de magias, rituais e a compreensão de fenômenos ocultos. Baixa Sabedoria não implica ignorância total, mas limita profundidade e alcance intelectual.

% =========================
% SEN
% =========================
\originlink{att:sense}
\noindent
\textbf{\baskerfamily\large SEN — Sense (Sentido)}\\
A percepção aguçada do mundo. SEN governa testes de observação, intuição e detecção de ameaças ocultas. Personagens atentos percebem o que outros ignoram — sons, mentiras, presenças invisíveis.


% =========================
% DIV
% =========================
\originlink{att:divinity}
\noindent
\hypertarget{target:divinity}{\textbf{\baskerfamily\large DIV — Divinity (Divindade)}}\\
A ligação do jogador com o Pináculo. DIV mede a proximidade do personagem com entidades superiores, deuses e forças cósmicas para permitir canalizar milagres, bênçãos e também maldições. DIV é o único atributo que tem limite máximo sendo ele 100 pontos.

\section*{E Afinal, o que é Divindade em Overgrown?}

Divindade representa o nível de progressão do personagem dentro do sistema de Overgrown. A cada nível de Divindade alcançado, o jogador recebe 5 pontos para serem distribuídos na árvore de talentos. Ao atingir 100 pontos de Divindade, o personagem alcança o limite máximo de conhecimento que uma mente humana é capaz de suportar do Pináculo, tornando-se o mais próximo possível de uma entidade divina.

Uma Divindade elevada não impõe, obrigatoriamente, mudanças no comportamento do personagem. Contudo, para fins de interpretação (\textit{roleplay}), seus efeitos narrativos podem ser definidos de comum acordo entre o mestre e o jogador.

No universo de Overgrown, é possível adquirir divindade de várias formas, seja entrando em contato com vestígios restantes de Energia Primordial, estudando conhecimentos divinos ou por meio de Combate. Todas as criaturas e seres tem vestígios de divindades, essa energia não se dissipa, e matar algo que contem divindade resulta na distribuição igualitária dela para os jogadors que participaram do combate. Considere a tabela para representações numéricas de cada nível de divindade:

\begin{center}
\renewcommand{\arraystretch}{1.3}
\begin{tabular}{@{} cc @{}}
\toprule
\rowcolor{gray!15} \textbf{Intervalo de Divindade} & \textbf{Pontos por Nível} \\ 
\midrule
01 -- 10 & 10 pts \\
11 -- 20 & 20 pts \\
21 -- 30 & 40 pts \\
31 -- 40 & 80 pts \\
41 -- 50 & 160 pts \\
\midrule
51 -- 60 & 320 pts \\
61 -- 70 & 640 pts \\
71 -- 80 & 1.280 pts \\
81 -- 90 & 2.560 pts \\
91 -- 100 & 5.120 pts \\
\bottomrule
\end{tabular}
\end{center}



\section*{Estatísticas dos jogadores}

Além dos atributos, jogadores tem algum recursos que são usado durante sua aventura: HP, IEP, Armadura e Esquiva.

HP, Health Points ou Pontos de Vida, representam a vida total do jogador, o quanto de dano ele consegue aguentar e é diretamente atrelado a sua \link{att:fortitude}{Fortitude}, cada ponto de Fortitude garante ao jogador 5 de vida. Um jogador com 10 de Fortitude terá 70 de vida, excluindo bônus e reduções de vida provenientes da árvore de talentos e de itens equipados.

IEP, Impossible Energy Points ou Pontos de Energia Impossível \link{res:null}{(Null)}, representa a capacidade do jogador de fazer feitos impossíveis, como usar Magias e Maldições. O IEP é diretamente ligado ao \link{att:wisdom}{Wisdom} do jogador, para cada ponto de Wisdom garante 5 pontos de IEP ao jogador. Um jogador com 10 de Wisdom terá 70 de IEP, excluindo bônus e reduções de IEP provenientes da árvore de talentos e de itens equipados.

Armadura reduz dano de ataques diretos. É diretamente baseada também na \link{att:fortitude}{Fortitude} do jogador. A cada 2 Pontos de Fortitude o jogador ganha 1 ponto de Armadura, e logo 1 de redução de ataques diretos. Um jogador com 20 pontos de Fortitude terá 10 de redução de dano, excluindo bônus e reduções provenientes da da árvore de talentos e de itens equipados.

\begin{rpgboxtips}{Bloqueio da Armadura}
A redução de dano proveniente da Armadura só é aplicada se o alvo do ataque decidir bloquear, em caso de esquivas e falha, apenas redução de danos nativas são contabilizadas. Bloqueio não conta como reação.
\end{rpgboxtips}

Esquiva define a proficiência em desviar do jogador. Ela é baseada diretamente ao atributo de \link{att:grace}{Grace}. A cada 2 Pontos de Grace, o jogador ganha 1 ponto de Esquiva. Todos os jogadores tem Esquiva base em 10. Um jogador com 20 de Grace, tem 20 de esquiva (10 base + 10 dos Pontos de Atributo), excluindo bônus e reduções provenientes da árvore de talentos e de itens equipados. Caso o rolamento de ataque seja menor que a esquiva do jogador, todo o dano é negado.

\begin{rpgboxtips}{Esquiva em Combate}
Esquivar de um ataque é considerado uma reação, alvos que esquivam perdem a reação na rodada. Falhar na esquiva ainda gasta sua reação.
\end{rpgboxtips}